\pdfoutput=1
\documentclass[11pt]{amsart}
\usepackage{amsmath,amssymb,amsthm,mathtools}
\usepackage[margin=1.2in]{geometry}
\usepackage{booktabs,array,adjustbox,needspace}
\usepackage[colorlinks=true,linkcolor=blue,citecolor=blue,urlcolor=blue]{hyperref}
\hypersetup{pdfauthor={Bernd Johannes Wuebben},pdftitle={Partial smoothings of Calabi--Yau threefolds with degree-six del Pezzo cone points}}
\setlength{\emergencystretch}{2em}
\newtheorem{theorem}{Theorem}[section]
\newtheorem{proposition}[theorem]{Proposition}
\newtheorem{lemma}[theorem]{Lemma}
\newtheorem{corollary}[theorem]{Corollary}
\theoremstyle{definition}\newtheorem{example}[theorem]{Example}
\theoremstyle{remark}\newtheorem{remark}[theorem]{Remark}
\newcommand{\CC}{\mathbb C}
\newcommand{\QQ}{\mathbb Q}
\newcommand{\ZZ}{\mathbb Z}
\newcommand{\PP}{\mathbb P}
\newcommand{\cO}{\mathcal O}
\DeclareMathOperator{\Def}{Def}
\title[Partial smoothings of Calabi--Yau threefolds]{Partial smoothings of Calabi--Yau threefolds\\ with degree-six del Pezzo cone points}
\author{Bernd Johannes Wuebben}
\date{26 November 2026}
\subjclass[2020]{14J32 (primary); 14B07, 14D15, 14E15 (secondary)}
\keywords{Calabi--Yau threefold, partial smoothing, deformation of singularities, del Pezzo surface, vanishing period}
\begin{document}
\begin{abstract}
Let $X$ be a projective Calabi--Yau threefold with ordinary double points
and anticanonical cones over smooth del Pezzo surfaces of degree six.
Assume that $X$ has a smooth projective crepant resolution which is small
at the nodes and blows up the cone vertices. Each cone has two local
smoothing components. For every choice of these components we prove that
their scheme-theoretic inverse image in the semiuniversal deformation
base of $X$ is smooth. If no global first-order deformation on this base
moves a given cone, its local classifying map is constant. Blowing up
precisely these vertices identifies the completed selected deformation
bases before and after the blow-up. On the resulting smooth base the
vanishing periods of a simultaneous nodal modification give a submersion
onto the space of relations among its exceptional curves. The zero sets of the coordinate
functions determine exactly which singularities can persist together on
a projective deformation. We describe the integral degree-two cohomology
of the regular locus of each partial smoothing and compute the restrictions
to the remaining links from divisor classes on the resolution.
\end{abstract}
\maketitle
\section{Introduction}
A cone over a degree-six del Pezzo surface admits two ways of smoothing.
Its reduced local deformation base is a line and a plane meeting at the
origin. A compact threefold can have many such points, with a global
deformation moving several at once. Fix either the line or the plane at
each vertex. What happens if the global first-order deformation map
vanishes at one of the chosen components? A vanishing derivative alone
would allow the cone to move at higher order. Even when all tangents
integrate, an arc tangent to a local discriminant can leave it.

We address both questions for a connected normal projective complex
threefold $X$ with $\omega_X\simeq\cO_X$ and $H^1(X,\cO_X)=0$.
Its singularities are ordinary double points and germs analytically
isomorphic to the anticanonical cone over the smooth degree-six del
Pezzo surface $E=\operatorname{Bl}_3\PP^2$. We assume a smooth projective
crepant resolution $Y\to X$, small at each node and given by the vertex
blow-up at each cone. Choose markings of the local germs, compatible
small resolutions in the local nodal modifications described below, and
analytic classifying maps from a semiuniversal base of $X$ to
semiuniversal bases of these germs. For a choice $S$ of the local
components, let $B=\Def^S(X)$ be the scheme-theoretic inverse image
of their product. At the ordinary nodes the full local bases are retained.

We recall the homological description of its tangent space.
A node contributes one signed exceptional curve. Over a cone the six
boundary curves of $E$ have self-intersection $-1$ and form a cycle.
The local line contributes the difference of two opposite boundary
curves. The plane contributes the two independent differences of
three alternating boundary curves. These choices mark the $A_1$ and
$A_2$ summands orthogonal to $K_E$. Denote the selected rational summand
at each singularity $p$ by $R_{p,S}$, including the exceptional line
at an original node. The inclusion of these curves in $Y$ gives
\[
 K_S=\ker\!\left(\bigoplus_p R_{p,S}\longrightarrow H_2(Y;\QQ)\right).
\]
Contraction with one global volume form identifies $K_S\otimes\CC$
with the image of $T_0B$ in the selected local tangent spaces
\cite[Lemma~4.4 and Theorem~4.5]{W5}. The earlier smoothness theorem
assumed that $K_S$ projects nontrivially at every cone. The first result
of this paper, Theorem~\ref{thm:regular}, removes this hypothesis for
the singularities considered here.

The step that permits its removal is to blow up the vertices at which
the chosen projection is zero. Over a small Stein neighborhood of such
a vertex, the resulting smooth neighborhood of $E$ is rigid. Its
contraction in any formal deformation of the blow-up is therefore locally
trivial. At the remaining vertices the earlier smoothness theorem
applies. Blow-down induces isomorphisms both on locally trivial tangents
and on the relation spaces. We obtain an isomorphism on the completed
local rings. This detects the scheme structure even where the images
of global first-order deformations miss an entire local component.

To determine partial smoothings we use the simultaneous modifications
over the local line and plane. Their central fibres replace a cone by
two or three nodes, respectively, and the new nodes have marked
exceptional curves on $Y$. Write $C_1,\ldots,C_n$ for these and the
original exceptional lines. The degree on each exceptional surface
forces the coefficients in its block to sum to zero. Consequently the
complex space of relations $R$ among the $C_i$ is canonically
$K_S\otimes\CC$ in these markings. A common holomorphic three-form
gives a period for each node of the modification. Stokes' theorem and
the tangent-space identification make this period map a submersion
onto $R$ (Theorem~\ref{thm:marked-periods}). In particular all coordinate
sections of $R$, including those on which a whole block vanishes, have
their actual geometric meaning on the selected analytic base. A nonzero
block on the plane with one zero period gives a residual node; an
identically zero block on an arc leaves the original cone.

For ordinary double points the description by relations among exceptional
curves goes back to Friedman \cite{Fried}. Lee, Lin and Wang identify
vanishing periods as coordinates transverse to the equisingular locus in
a projective conifold transition \cite[Lemma~3.13 and Proposition~3.16]{LLW}.
Namikawa studies strata obtained from small projective modifications of
terminal Calabi--Yau threefolds and compares deformation spaces of
terminal small modifications \cite[Proposition~2.3 and Theorems~2.5,~2.7]{Nam}.
The degree-six cone has a crepant exceptional divisor. Its local line
and plane are the two components of Altmann's versal toric deformation
\cite{Alt}; Gross discusses them in his analysis of primitive
contractions \cite[\S5]{Gross}. The full-smoothing criterion of
\cite[Theorem~4.7]{W5} asks whether a global relation avoids every
local discriminant. Here the nodal intermediate may itself be
nonprojective or nonsmoothable. We work on the selected deformation
base of $X$ and retain the marked zero periods as well as the nonzero
ones.

Finally, Proposition~\ref{prop:marked-generators} determines the
degree-two classes on the regular locus of a partially smoothed fibre
from their evaluations on the marked curves. In a fluxless reduction
of M-theory these data identify the surviving bulk gauge generators,
including their action on the remaining links. Two compact hypersurface
examples illustrate the difference between the rank of this group and
its measured charges. They also give choices of components for which
the omitted nonzero-projection hypothesis fails.

\section{Deformations at cone points with zero projection}
\label{sec:rigidity}

Fix a component $S_p$ of the local deformation base at each cone point
$p$, and let $B=\Def^S(X)$ be the inverse image of their product under
the local classifying maps. We use the scheme structure on this inverse
image. Write $F$ for the set of cone points at which the projection of
$K_S$ to the selected local summand is zero. We first study deformations
of the vertex blow-up at a point of $F$.

\begin{lemma}\label{lem:blowup-rigidity}
Let $U$ be a sufficiently small Stein representative of the
anticanonical cone over a smooth del Pezzo surface $E$ of degree six.
Let $f:V\to U$ be the blow-up of its vertex. Then
\[
 f_*\Theta_V=\Theta_U,\qquad R^1 f_*\Theta_V=0.
\]
Every marked deformation of $V$ over a local Artin $\CC$-algebra is
trivial. Across a small extension the trivialization can be chosen to
extend any specified trivialization of the preceding order.
\end{lemma}

\begin{proof}
The singular locus of $U$ consists of the vertex. The local flow of a
vector field on $U$ therefore preserves the vertex ideal and lifts to
its blow-up. In the opposite direction a vector field descends on the
punctured germ and extends across the vertex by reflexivity of
$\Theta_U$. This proves the direct-image assertion.

Near the exceptional divisor $V$ is a neighborhood of the zero section in
the total space of $K_E$. If $I$ is its ideal, the graded quotients of
$\Theta_V$ for the $I$-adic filtration are
\[
 \bigl(\Theta_E\oplus K_E\bigr)\otimes\mathcal O_E(-kK_E),
 \qquad k\geq 0.
\]
Their first cohomology vanishes, as in the local calculation
in~\cite[proof of Lemma~5.2, (47)--(48)]{W5}. Indeed,
$H^1(E,\mathcal O_E((1-k)K_E))=0$ by rationality, Serre duality and
Kodaira vanishing. For the tangent summand, the toric Euler sequence
has middle term $\bigoplus_{j=1}^6\mathcal O_E(C_j)$, where the $C_j$
are the six boundary curves, and left term $\mathcal O_E^{\oplus4}$.
When $k=0$, the restriction $\mathcal O_{C_j}(C_j)$ has degree $-1$,
giving $H^1(E,\mathcal O_E(C_j))=0$. When $k\geq1$, the divisor
$C_j-kK_E$ has nonnegative degree on each of the six generators of the
cone of curves, so toric vanishing applies. The last term needed in the
cohomology sequence is $H^2(E,\mathcal O_E(-kK_E))=0$.
Induction on the thickenings of $E$ and the theorem on formal functions
give $R^1f_*\Theta_V=0$. Here the powers of $I$ are cofinal with the
pullbacks of the powers of the vertex ideal.

Cartan's theorem~B and Leray now give $H^1(V,\Theta_V)=0$. Consider a
small extension of local Artin algebras with kernel $J$. Once a
trivialization over its quotient is fixed, the differences of local
trivializations lifting it define a class in
$H^1(V,\Theta_V)\otimes J$. The local lifts exist since $V$ is smooth.
The class vanishes, so the lifts can be made to agree. Induction from
the given central marking proves the assertion.
\end{proof}

We shall use the following consequence of semiuniversality. Let $(G,0)$
be a semiuniversal base of an isolated analytic singularity $(U,p)$,
and let $D$ be a formal power-series ring over $\CC$. If a map
$\operatorname{Spf}D\to\widehat G$ induces a trivial marked local
deformation at every Artin quotient, that map is constant. To see this,
choose a generic formal line through the origin on which a first nonzero
term survives, and call its order $m$. The map modulo $t^{m+1}$ factors
through the first infinitesimal neighborhood of $0\in G$, with tangent
coefficient in $(t^m)/(t^{m+1})$. The Kodaira--Spencer isomorphism
identifies this coefficient with a nonzero element of
$T^1_{U,p}\otimes (t^m)/(t^{m+1})$, using the product trivialization
modulo $t^m$. A change of lower-order trivialization cannot erase this
class. In fact every marked automorphism of the product modulo $t^m$
lifts to the product modulo $t^{m+1}$. Its logarithm is a derivation
with coefficients in $(t)$: the finite binomial formula for powers of a
unipotent automorphism is multiplicative at every integer exponent,
hence at every complex exponent, and differentiation at zero gives
the Leibniz rule. Lift the coefficients of this derivation and
exponentiate. Any automorphism of the central germ lifts constantly.
Thus a trivial deformation modulo $t^{m+1}$ would have zero class for
the specified product identification, a contradiction.

\begin{theorem}\label{thm:regular}
For every choice of the local components $S_p$, the analytic germ
$\Def^S(X)$ is smooth of dimension
\[
 h^1(X,\Theta_X)+\dim_{\QQ} K_S.
\]
At every $p\in F$ the local classifying map is constant on this entire
germ. Let $X_F\to X$ be the blow-up of the vertices in $F$, and let
$S'$ denote the induced choices at the remaining cone points.
Blow-down induces an isomorphism of the completions of
$\Def^{S'}(X_F)$ and $\Def^S(X)$.
\end{theorem}

\begin{proof}
The blow-up $\pi:X_F\to X$ is projective and crepant. It replaces each
vertex in $F$ by the smooth zero section of its canonical bundle and
is an isomorphism near the remaining singularities. Rationality of
these cone singularities gives $R\pi_*\mathcal O_{X_F}=\mathcal O_X$.
In particular $X_F$ is projective with trivial dualizing sheaf and
$H^1(X_F,\mathcal O_{X_F})=0$. The same smooth threefold $Y$ resolves
$X$ and $X_F$. Its marked curves at the remaining singularities are
unchanged. Extension of a tuple by zero at $F$ identifies
$K_{S'}(X_F)$ with $K_S$, since every relation in $K_S$ already
vanishes there. In particular every remaining selected cone summand
has nonzero projection.

Write $B'=\Def^{S'}(X_F)$. The selected-base theorem
\cite[Theorem~4.5]{W5} makes $B'$ smooth. Let
$D=\widehat{\mathcal O}_{B',0}$. Analytic blow-down over Artin
algebras gives a formal morphism from $\operatorname{Spf}D$ to
$\widehat{\Def}(X)$ compatible with restriction to open sets and with
base change \cite[Lemmas~4.2--4.3]{W5}. At each vertex in $F$,
Lemma~\ref{lem:blowup-rigidity} trivializes its source over the
inverse image of a Stein neighborhood, to every order. Pushing forward
the product structure sheaf gives the trivial deformation of that
neighborhood. The preceding consequence of semiuniversality shows
that all coordinates of its fixed local classifying map therefore
vanish in $D$.

At every other cone the morphism $\pi$ is a local isomorphism. The
differential of blow-down on its local components agrees with the
original restriction, using the pullback of one global volume form
for the homological identifications. The selected tangent sequence
\cite[Lemma~4.4]{W5} for $X_F$ and the nonzero projection at this cone
show that at least one coordinate of the chosen component has nonzero
linear term in $D$. The completed local deformation ring of the cone is
$\CC[[s,u,v]]/(su,sv)$, with line $(u,v)=0$ and plane $s=0$.
Since $D$ is a domain, the complementary coordinates vanish. Thus the
formal morphism factors through $\widehat B$ at every cone. Nodes
impose no further choice. Denote the factored morphism by $\beta$.

The map $d\beta$ identifies the two selected tangent spaces. Indeed,
Lemma~\ref{lem:blowup-rigidity} and the local isomorphism away from $F$
give $\pi_*\Theta_{X_F}=\Theta_X$ and $R^1\pi_*\Theta_{X_F}=0$.
Leray identifies $H^1(X_F,\Theta_{X_F})$ with $H^1(X,\Theta_X)$.
On locally trivial deformations this is the differential of blow-down:
a vector-field cocycle on overlaps descends by the direct image of
its vector fields. On the local quotient the differential extends the
tuple by zero at $F$, hence gives the identification of relation
spaces above. For each of the pairs $(T,S_T)=(X_F,S')$ and $(X,S)$
we have the exact sequence
\[
 0\longrightarrow H^1(T,\Theta_T)
 \longrightarrow T_0\Def^{S_T}(T)
 \longrightarrow K_{S_T}(T)\otimes_{\QQ}\CC\longrightarrow0.
\]
The isomorphisms on its two ends give the assertion about $d\beta$.

Write $\widehat{\mathcal O}_{B,0}=\CC[[x_1,\ldots,x_r]]/I$
minimally. The images of the $x_i$ in the regular complete local ring
$D$ form a basis of its cotangent space. They form a system of formal
coordinates, so the composite $\CC[[x_1,\ldots,x_r]]\to D$ is an
isomorphism. It follows that $I=0$ and that $\beta$ is an isomorphism.
Regularity of the completion proves smoothness of $B$, and the
tangent sequence gives its dimension. At any $p\in F$ the local
classifying coordinates vanish after pullback to $D$. They therefore
vanish in $\widehat{\mathcal O}_{B,0}$ and in its analytic local
ring. This proves constancy on $B$.
\end{proof}

\section{Periods and the discriminant on a selected deformation space}
\label{sec:marked-periods}

Write $B=\Def^S(X)$. At each degree-six vertex take the simultaneous
incidence modification on the chosen local component. Its central fibre
has two ordinary double points on the line component and three on the
plane component. All the original ordinary double points are retained.
Write $Z\to X$ for this modification on the central fibre and
$C_1,\ldots,C_n$ for the exceptional curves of the compatible small
resolution $Y\to Z$. The index of each curve records its original node
or its position over a marked cone. Orient the curves over a cone by
their complex orientations. At an original node a change of small
resolution changes the sign of its marked class. Put
\[
 R=\ker\!\left(\CC^n\longrightarrow H_2(Y;\CC),\quad
                         w\longmapsto\sum_iw_i[C_i]\right).
\]
The support of $w$ is the set of indices with $w_i\ne0$.

\begin{theorem}\label{thm:marked-periods}
The marked nodal periods define a holomorphic submersion of germs
$P\colon(B,0)\to(R,0)$. There are local analytic coordinates on $B$
in which $P$ is projection onto $R$. For each $i$ the ordinary double
point of the simultaneous modification persists exactly on
$P_i^{-1}(0)$, with its scheme structure as the pullback of the local
nodal discriminant. In particular every support occurring in $R$ occurs
on a projective deformation of $X$ inducing the selected profile.
\end{theorem}

\begin{proof}
Restrict the analytic family of $X$ over $B$ to representatives of the
marked germs. Pull back the simultaneous families of
\cite[\S3.4 and Proposition~3.8]{W5} along their classifying maps.
The modifications are proper, are flat over the parameter line or
plane, and are isomorphisms away from the vertex over the local
parameter origin. Thus the pullbacks glue to the identity on the
complement of smaller vertex neighborhoods. Properness is local on
the target, and flatness over $B$ follows from the two types of
charts. This gives a proper flat analytic family $g\colon\mathcal Z\to B$
with central fibre $Z$. A block whose local parameter is zero has the
central modification over that vertex, even on a fibre where another
vertex has moved. At a nonzero local parameter the map to the
corresponding part of the fibre of $X$ is an isomorphism, including
on the plane discriminant.

We use one relative volume form for all the nodes. Indeed, $Z$ is
connected and Gorenstein with $\omega_Z\simeq\cO_Z$ and
$H^1(Z,\cO_Z)=0$ by \cite[Proposition~4.1]{W5}. Serre duality
gives $H^2(Z,\cO_Z)=0$. After shrinking $B$, the relative dualizing
sheaf is invertible and commutes with base change. Semicontinuity
gives $h^0(\cO_{Z_b})=1$ and $h^1(\cO_{Z_b})=h^2(\cO_{Z_b})=0$.
Constancy of the holomorphic Euler characteristic in the proper flat
family gives $h^3(\cO_{Z_b})=1$. By duality and Grauert's base change
theorem a generator of $H^0(Z,\omega_Z)$ extends to a relative section
$\Omega$. Its zero locus misses the central fibre, so properness lets
us shrink $B$ until $\Omega$ generates $\omega_{\mathcal Z/B}$.

Remove mutually disjoint Milnor neighborhoods of the $n$ nodes of
$Z$. The compact exterior $W$ and its boundary are differentiably
locally trivial over $B$, by the proper submersion theorem with
boundary. The boundary component $L_i$ is diffeomorphic to
$S^2\times S^3$. Let $\gamma_i$ generate $H_3(L_i;\ZZ)$, oriented
so that the boundary of the Poincar\'e--Lefschetz dual on $W$ of
$D|_W$, for $D\in H^2(Y;\ZZ)$, has $i$th component
$\langle D,C_i\rangle\gamma_i$. This convention is possible since
restriction from a small-resolution neighborhood to the link is an
isomorphism on $H^2$. Transport the boundary cycles and set
$P_i(b)=\int_{\gamma_i}\Omega_b$. On the exterior $\Omega_b$ is
closed. Stokes' theorem on the transported relative four-cycle gives
\[
              \sum_i\langle D,C_i\rangle P_i(b)=0 .
\]
Integral degree-two classes span $H^2(Y;\CC)$ over $\CC$. The
evaluation pairing therefore gives $P(b)\in R$ on every fibre.

In a nodal chart write $\sum_{j=1}^4 x_j^2=\tau_i(b)$. The relative
form is $u(x,b)$ times the residue form, for $u$ a holomorphic unit.
When $\tau_i(b)\ne0$, inclusion of the link into the local Milnor
fibre induces an isomorphism on $H_3$: the fibre is a disc cotangent
bundle of $S^3$, with trivial sphere bundle as boundary. Up to the
chosen sign, $\gamma_i$ thus represents the sphere
$\sqrt{\tau_i(b)}\,S^3$. Scaling the residue form shows that its
integral is a fixed nonzero constant times $\tau_i(b)$ times the
spherical average of $u(\sqrt{\tau_i(b)}\,x,b)$. The odd Taylor
terms have zero average and the even terms converge normally on
a smaller chart. The average is holomorphic in $b$ and has nonzero
constant term $u(0,0)$. When $\tau_i(b)=0$, the cycle bounds in
the local small resolution and the form extends there by crepancy.
Consequently, also for an identically persistent node, we have
\begin{equation}\label{eq:period-unit}
               P_i(b)=\tau_i(b)v_i(b),\qquad v_i(0)\ne0
\end{equation}
with $v_i$ holomorphic. The units in this calculation come from the
same global form. Rescaling that form multiplies the whole vector
$P$ by a single unit on $B$.

Intersecting a relation with a vertex divisor $E_p$ gives
$\sum_{i\text{ over }p}w_i=0$: each of its marked curves has
$E_p$-degree $-1$, and all the other marked curves have degree zero.
The line block is consequently $(\lambda,-\lambda)$, and the plane
block can be written $(-a,b,a-b)$. The curve identifications of
\cite[Proposition~3.7 and Lemma~4.4]{W5} identify $R$ in these
coordinates with $K_S\otimes_{\QQ}\CC$. The local-to-global
tangent sequence in that lemma gives the image of the differential
of the selected local classifying map, of dimension $\dim R$.
Its subsequent differential to the tuple of nodal parameters is
injective on the local target: it is the identity at an original
node, records a nonzero multiple of the line parameter at a line
vertex, and on the plane records the three differences of the three
diagonal entries modulo simultaneous translation. These differences
separate tangent vectors in the plane. By \eqref{eq:period-unit},
multiplication by the diagonal matrix with entries $v_i(0)$ leaves
this rank unchanged. Since $P$ takes values in $R$, its differential
at zero is surjective.

Theorem~\ref{thm:regular} and the holomorphic submersion theorem
now give the asserted coordinates. Given $w\in R$, take a
local section through the origin and lift $t\mapsto tw$. Equation
\eqref{eq:period-unit} puts this entire disc in every prescribed
zero divisor and off every other one for $t\ne0$. A zero block is the
branch origin, since the nodal parameters jointly separate the
points of each local line or plane. A nonzero plane block with one
zero coordinate leaves precisely one ordinary double point. Outside
the marked neighborhoods the family is smooth after shrinking.

Let $f\colon\mathcal X\to\Delta$ denote the resulting disc family
of $X$, and fix an ample bundle $L$ on $X$. Topological proper base
change identifies $H^2(X;\ZZ)$ with the direct limit of
$H^2(f^{-1}(U);\ZZ)$ over discs $U$ about zero. Thus $c_1(L)$
extends after shrinking. Serre duality gives $H^2(X,\cO_X)=0$.
By semicontinuity, Grauert's theorem and the Stein property of the
disc, $H^2(\mathcal X,\cO_{\mathcal X})=0$ there. The exponential
sequence lifts the extended class to a line bundle, which restricts
to $L$ since $H^1(X,\cO_X)=0$. Openness of ampleness gives a
projective family even when $Z$ is nonprojective.
\end{proof}

For a projective conifold transition this arrangement is the one
described by Lee, Lin and Wang \cite[Lemma~3.13 and
Proposition~3.16]{LLW}. In the present construction $Z$ may itself
have a node which persists on the whole base. Such an index contributes
the zero functional on $R$. Two nonzero proportional coordinate
functionals give the same hyperplane, with both node markings retained.
The submersion concerns periods. The local equation parameters in
\eqref{eq:period-unit} may differ from them by different units.

\section{Cohomology of the regular locus and the marked torus}
\label{sec:marked-generators}

Let $I$ be a support in $R$. Denote by $F(I)$ the set of degree-six
vertices for which every index over the vertex lies outside $I$.
Put $H=H^2(Y;\ZZ)$, write $[E_p]\in H$ for the Poincar\'e dual
of the vertex divisor, and let
\[
 A_I=\{D\in H:\langle D,C_i\rangle=0\ \text{for all }i\in I\}.
\]
For a nearby fibre $X_b$ with $\operatorname{supp}P(b)=I$, the
following statement uses the markings transported from the central
fibre. Write $X_b^{\rm reg}$ for its smooth locus.

\begin{proposition}\label{prop:marked-generators}
There is an identification
\[
 H^2(X_b^{\rm reg};\ZZ)
       \simeq A_I\Big/\sum_{p\in F(I)}\ZZ[E_p].
\]
It identifies this group with a subgroup of
$H^2(X^{\rm reg};\ZZ)=H/\sum_p\ZZ[E_p]$. A class with lift $D\in H$
belongs to the image precisely when it has degree zero on every
original node in $I$ and its degrees on the curves in $I$ over each
one vertex are equal. For a representative $D\in A_I$, the restriction
at a residual ordinary double point is its degree on the marked
exceptional line. At a persistent cone $p$ it is
\[
 [D|_{E_p}]\in H^2(E_p;\ZZ)/\ZZ K_{E_p}\simeq H^2(L_p;\ZZ),
\]
where $L_p$ denotes the link of that cone.
\end{proposition}

\begin{proof}
Use the coordinates of Theorem~\ref{thm:marked-periods} to join the
origin to $b$ by an arc along which the support is $I$ away from the
origin. On a block in $F(I)$ the incidence family is the product of
the central model with that arc. Fill the exterior $W$ at the
indices in $I$ by the Milnor fibres, and at the remaining indices
by the compatible local small resolutions. Write $M_I$ for the
resulting smooth compact manifold. The filling which uses small
resolutions at all indices gives $Y$.

For either filling, each $H^1(L_i;\ZZ)$ is zero. On $H^2$ the
restriction of a small-resolution neighborhood to its boundary is an
isomorphism, and the Milnor fibre has $H^2=0$. Mayer--Vietoris first
identifies $H^2(Y;\ZZ)$ with $H^2(W;\ZZ)$ and then identifies the
image of $H^2(M_I;\ZZ)$ there with $A_I$. A vertex in $F(I)$
reconstructs the original smooth surface $E_p\subset M_I$, with
the marking of its divisor class preserved in the common exterior.
Remove these pairwise disjoint surfaces. By the Thom sequence,
the degree-two cohomology of the complement is obtained by quotienting
by their divisor classes, since $H^1(E_p;\ZZ)=0$. The exceptional
curves at all other retained nodes are disjoint from the removed
surfaces. Removing them has no effect in degree two, by the Thom
isomorphism in real codimension four. This complement is the smooth
locus of $X_b$. The case $I=\varnothing$ gives the central
identification in the statement.

The resulting map to $H^2(X^{\rm reg};\ZZ)$ is induced by
$A_I\hookrightarrow H$. A combination $\sum_pc_p[E_p]$ which
annihilates the curves in $I$ has $c_p=0$ for each vertex with an
index in $I$, as one sees by evaluating on any such curve. This
gives injectivity on the stated quotient. Changing a lift by
$c_p[E_p]$ subtracts the integer $c_p$ from each degree over $p$
and has no effect at other vertices or original nodes. The degrees
in a nonempty block can all be made zero precisely when they are
equal. At an ordinary double point restriction to the link agrees
with the degree on the exceptional line by the same filling calculation.
At a persistent cone the link is the circle bundle of $K_{E_p}$.
The Gysin sequence gives the stated quotient since $H^1(E_p;\ZZ)=0$;
adding $[E_p]$ to the lift adds $K_{E_p}$ to the restriction by adjunction.
This also describes the link classes outside the selected root summand.
\end{proof}

For the rigid abelian gauging, one can see both the supports and
the stabilizers directly. Work at zero triplet moment-map level with
just the marked fields. Use one charged hypermultiplet at each
original node and the presentation by two or three homogeneous
hypermultiplets at a chosen cone component, divided there by the
diagonal circle \cite[\S\S2--3]{Higgs}. The three moment maps for
one homogeneous hypermultiplet $(q_i,\widetilde q_i)$ comprise its
complex product $q_i\widetilde q_i$ and its real moment
$(|q_i|^2-|\widetilde q_i|^2)/2$. Their joint vanishing is equivalent
to $(q_i,\widetilde q_i)=0$. A common rotation of the three moment
maps makes the complex projection nonzero on all the nonzero
triplets of a given solution: the discarded real axis need only
avoid finitely many lines. The resulting complex products satisfy
the linear charge relations and give a vector of $R$ with exactly
that support. Conversely, for $w\in R$ take each product equal to
$w_i$ and both absolute values equal to $\sqrt{|w_i|}$. This solves
all three equations, including each diagonal-circle equation.
Together with Theorem~\ref{thm:marked-periods}, this gives equality
of the sets of supports realized on the two sides.

For the connected bulk torus take the cocharacter lattice
$H^2(X^{\rm reg};\ZZ)$ modulo torsion, with the original marked
degrees. Its action fixes the class of a configuration precisely
when a diagonal-circle action at each cone compensates the bulk
action on every nonzero homogeneous hypermultiplet there. On the
Lie algebra this requires equal degrees on the participating curves
over that cone and zero degree at each participating original node.
For a cocharacter with integral degrees the compensating common
degree is integral as well. Proposition~\ref{prop:marked-generators}
thus gives the marked lattice of the identity component of the
stabilizer, including spectator directions. An empty block leaves
its diagonal circle redundant, as accounted for by the quotient
by $[E_p]$. Finite components of a stabilizer depend in addition
on saturation of the subgroup of the character lattice generated
by the active weights.

With a conical Ricci-flat metric on the regular locus, its
degree-two $L^2$ cohomology identifies with middle-perversity
intersection cohomology \cite{Cheeger}. At an isolated point of
real codimension six, the middle-perversity truncation retains
the link cohomology through degree two. Thus
$IH^2(X_b;\mathbb R)=H^2(X_b^{\rm reg};\mathbb R)$. In a fluxless
three-form reduction these classes give the linearized normalizable
bulk vector fields. The local conicality criterion for bounded-potential
Ricci-flat metrics at isolated Calabi--Yau cones with smooth link is
given in \cite[Theorem~1.4]{Zhang}. The cohomological calculation
identifies the marked vector space on the geometric fibre. A
microscopic identification of the complete Higgs vacuum at finite
Planck mass would require further information about the localized
interacting theory.

\section{Two compact hypersurfaces}\label{sec:examples}

Let $X_-$ be the anticanonical threefold for the fan polytope
$\Delta_{25}$ of \cite[\S5]{W5}. For the second example take the
four-dimensional fan polytope whose first four vertices are the
standard basis vectors and whose remaining vertices, in order, are the
columns of
\[
{\setlength{\arraycolsep}{1.2pt}\scriptscriptstyle
\left(\begin{array}{*{23}{r}}
-1&1&1&0&0&-1&-1&1&0&0&-1&0&-1&1&0&-1&1&-2&0&-1&-2&0&-1\\
1&-1&-1&0&0&1&1&-1&0&0&1&-1&0&-2&-1&0&-2&1&-1&0&1&-1&0\\
1&0&-1&1&-1&1&0&-1&0&-1&0&0&1&-1&-1&1&-1&1&0&-1&0&-1&-1\\
0&0&1&-1&1&-1&0&0&-1&0&-1&1&0&1&1&-1&0&-1&-1&0&-1&-1&-1
\end{array}\right).}
\]
This is polytope 528 in the census accompanying \cite{W5}. Number
its vertices from zero. Its four hexagonal faces, written in cyclic
order, are
\[
\begin{aligned}
 F_0&=(19,22,20,17,15,16),& F_1&=(18,15,16,21,24,23),\\
 F_2&=(25,26,23,18,17,20),& F_3&=(25,22,19,21,24,26).
\end{aligned}
\]
Take a general anticanonical hypersurface and subdivide at the
centres of $F_2,F_3$. Write $X_+$ for its proper transform. The two
threefolds $X_-,X_+$ each have thirty ordinary double points and two
anticanonical cones over del Pezzo surfaces of degree six. For $X_+$
this follows by taking the two-faces of the displayed convex hull.
The parallelograms are unit squares up to integral affine
equivalence, the other two-faces apart from the four hexagons are
unimodular triangles, and the dual edge of each parallelogram or
hexagon is primitive. The two centre subdivisions resolve the
corresponding cone points. Resolving the two remaining vertices and
refining projectively gives a common smooth crepant resolution
compatible with the local contractions used below. Its Picard rank is
$27$ for $X_+$ and $23$ for $X_-$. The only nonvertex boundary
lattice points are the hexagon centres, and the dual-edge lengths
give zero correction in the toric divisor count.

\subsection{Supports on the local components}
Write $L$ and $P$ for the line and plane of the local deformation
space. On both threefolds order the two components by $F_0,F_1$, using
the cyclic face markings of \cite[\S5]{W5} on $X_-$. Put the thirty
parallelograms in lexicographic order on their increasing lists of
vertices. The circuit has sign $+1$ at its least vertex and the
opposite vertex. On a marked hexagon let $B_0,\dots,B_5$ denote
the cyclic boundary curves of its vertex resolution. The curve
columns for $L$ are $B_0,B_3$ and those for $P$ are $B_0,B_2,B_4$.
The divisor degrees of $B_j$ are $-1$ at its vertex and the
centre, and $1$ at the two neighboring vertices. The remaining
degrees vanish. These rules give the entire curve matrix.\footnote{The
standalone SageMath supplement \texttt{check\_examples.sage}
reconstructs the faces from the two vertex lists and repeats the
integral calculations of this section.}

In Table~\ref{tab:maximal-support} the node indices run from $0$ to
$29$ and are followed by the two ordered groups of curve indices
over the marked cones. Write $r$ for the rank of all curve classes,
and $r_I$ for the rank of those on a support $I$ of greatest
cardinality. The indicated missing indices vanish in every relation.
The last column records the singularities of a fibre realizing $I$.
\begin{table}[ht]
\centering
\caption{Largest supports for four choices of local components.}
\label{tab:maximal-support}
\begin{tabular}{cccccccl}
\hline
& choice & curves & $r$ & $\dim R$ & $r_I$ & missing indices & singularities\\
\hline
$X_-$&$LP$&35&22&13&21&33&one node\\
$X_+$&$LP$&35&22&13&22&none&smooth\\
$X_+$&$LL$&34&23&11&20&17, 32, 33&node 17 and cone at $F_1$\\
$X_+$&$PP$&36&24&12&21&30, 31, 32&cone at $F_0$\\
\hline
\end{tabular}
\end{table}

Here are integral relations with the asserted supports. The entries
of each vector $d,p,q$ give the thirty old nodal coefficients in
three successive groups of ten.
\[
\begin{array}{c|*{10}{r}}
d&1&-2&1&1&1&-2&-1&1&-1&2\\
 &-2&2&1&1&-1&-1&-1&1&-1&1\\
 &1&1&-1&-2&2&-1&2&1&-2&2\\[2pt]
p&1&2&1&-2&-1&-2&2&-1&2&-1\\
 &1&-1&1&1&-1&2&-2&0&-1&1\\
 &2&-1&2&-1&1&-1&1&-1&-1&1\\[2pt]
q&-2&-1&1&1&-1&-1&2&1&-1&-1\\
 &-2&-2&2&-1&-1&1&1&-2&-1&1\\
 &1&-3&-1&2&-1&2&1&2&2&-2
\end{array}
\]
In table order the four relations are
\[
\begin{aligned}
 &(d,\ 1,-1,\ -1,0,1),&& (p+3q,\ 2,-2,\ 6,-9,3),\\
 &(p,\ 2,-2,\ 0,0),&& (q,\ 0,0,0,\ 2,-3,1).
\end{aligned}
\]
Each relation follows by summing the signed curve columns just
described. Row reduction gives the ranks in the table. Every nonzero
Smith invariant of each active matrix is one.

Write $D(j_1,\ldots,j_t)=\sum_{k=1}^tD_{j_k}$ for the sum of the
indicated prime divisors. The centre rays follow the vertices in face
order. On $X_-$ the divisor $D(2,3,13,14,15,16)$ has degree one on
curve $33$ and degree zero on all other $LP$ columns.
On $X_+$ with choice $LL$, the divisor $D(2,3,4,15,16)$ isolates
the coefficient of node $17$ after the sum-zero relation over
each cone is imposed. The divisor $D(3,6,8,15,17,18,23)$ similarly
isolates the line coefficient at $F_1$. Thus that cone and node
$17$ persist. The two divisor evaluations in the next paragraph
account for the vanished plane at $F_0$ in the last row.

Theorem~\ref{thm:marked-periods} realizes all four displayed
supports. For the two examples with a constant cone, their existence
can also be seen inside the $LP$ deformation space of $X_+$.
Both marked projections are nonzero there. To obtain the $LL$ fibre,
take the line at $F_0$ with coefficients $(2,-2)$ and keep the
plane at $F_1$ at its origin. To obtain the $PP$ fibre, keep the
line at $F_0$ at its origin and take the plane at $F_1$ with
coefficients $(2,-3,1)$. A locally trivial cone deformation belongs
to both components.

\subsection{Divisor classes at the persistent cone}
On each hexagon use the ordered root classes
\[
 \ell=B_0-B_3,\qquad A=B_4-B_0,\qquad B=B_2-B_4 .
\]
They form an integral basis of $K_E^\perp$ in the Picard group of
the degree-six del Pezzo surface. Consider the $PP$ example on
$X_+$. Put
\[
 D_A=-D(11,13,20,22,25,26),\qquad
 D_B=D(3,6,8,15,17,18,23,28).
\]
Both divisors annihilate the thirty old nodes and all three plane
curves over $F_1$. Their degrees on
$(\ell_0,A_0,B_0;\ell_1,A_1,B_1)$, where subscripts now specify
the face, are respectively
\[
 (0,1,0;0,0,0),\qquad (0,0,1;1,0,0).
\]
In particular they isolate the two independent plane coefficients
at $F_0$ when both components are planes. If $N$ denotes the
matrix of the thirty old curve rows, its rank is $20$. Appending
all six root rows increases its rank to $22$. Thus restrictions
of classes annihilating the old curves have rank at most two.
The two tuples above span a primitive sublattice of $\ZZ^6$,
so they give the whole integral image of the surviving classes in
the original marking. On the link of the persistent cone their
degrees span the two plane coordinates. The last three entries in
each tuple retain the marking at the other cone of the central fibre.

The regular locus of this fibre has second Betti number
$27-21-1=5$ by Proposition~\ref{prop:marked-generators}. There are
three further independent classes with zero marked degrees.
Let $H$ be the pullback of the anticanonical polarization on the
original toric fourfold, and let $E_2=D_{29}$ and $E_3=D_{30}$ be the
two exceptional surfaces retained in forming $X_+$. Together
$D_A,D_B,H,E_2,E_3$ give an integral basis modulo torsion. Indeed,
on the resolution one has $H^3=83$ and $H^2D_0=12$, the normalized
volumes of the polar polytope and its facet dual to $v_0$.
The integral class $7HD_0-H^2$ therefore pairs to one with $H$
and pairs to zero with every centre divisor. A boundary line on
each $E_2,E_3$ pairs to $-1$ with its own surface and to zero with
the other surface and $H$. These three functionals annihilate the
class of $E_0$ removed in Proposition~\ref{prop:marked-generators}.
They show that $H,E_2,E_3$ span a primitive sublattice of rank
three in the kernel of the marked degrees. The two displayed root
tuples then give the assertion.

\begin{thebibliography}{Wue26b}
\bibitem[Alt97]{Alt} K.~Altmann, \emph{The versal deformation of an isolated toric
Gorenstein singularity}, Invent. Math. \textbf{128} (1997), 443--479.
\bibitem[Che83]{Cheeger} J.~Cheeger, \emph{Spectral geometry of singular Riemannian spaces},
J. Differential Geom. \textbf{18} (1983), 575--657.
\bibitem[Fri86]{Fried} R.~Friedman, \emph{Simultaneous resolution of threefold double points},
Math. Ann. \textbf{274} (1986), 671--689.
\bibitem[Gro97]{Gross} M.~Gross, \emph{Deforming Calabi--Yau threefolds}, Math. Ann.
\textbf{308} (1997), 187--220; corrected version,
\href{https://arxiv.org/abs/alg-geom/9506022v2}{arXiv:alg-geom/9506022v2}.
\bibitem[LLW18]{LLW} Y.-P.~Lee, H.-W.~Lin and C.-L.~Wang, \emph{Towards $A+B$ theory in
conifold transitions for Calabi--Yau threefolds}, J. Differential Geom. \textbf{110}
(2018), 495--541.
\bibitem[Nam02]{Nam} Y.~Namikawa, \emph{Stratified local moduli of Calabi--Yau threefolds},
Topology \textbf{41} (2002), 1219--1237.
\bibitem[Wue26a]{W5} B.~J. Wuebben, \emph{Smoothing Calabi--Yau threefolds with nodes
and del Pezzo cone points}, \href{https://arxiv.org/abs/2610.35915}{arXiv:2610.35915}.
\bibitem[Wue26b]{Higgs} B.~J. Wuebben, \emph{Global constraints on Higgsing localized
five-dimensional sectors in compact M-theory vacua}, preprint (2026),
\href{https://doi.org/10.5281/zenodo.23643226}{doi:10.5281/zenodo.23643226}.
\bibitem[Zha25]{Zhang} J.~Zhang, \emph{On polynomial convergence to tangent cones for
singular K\"ahler--Einstein metrics}, J. Reine Angew. Math. \textbf{826} (2025), 271--302.
\end{thebibliography}
\end{document}
